\documentclass[aspectratio=169, 10pt]{beamer}

% --- Paquetes Fundamentales ---
\usepackage[utf8]{inputenc}
\usepackage[T1]{fontenc}
\usepackage[spanish]{babel}
\usepackage{amsmath, amssymb}
\usepackage{microtype}

% --- Matemáticas y Electrónica ---
\usepackage{siunitx}
\usepackage[american, RPvoltages]{circuitikz}

% --- Configuración de siunitx ---
\sisetup{output-decimal-marker = {,}, per-mode = symbol}

% --- Configuración de Colores e Identidad ---
\definecolor{ucbblue}{RGB}{0, 51, 102}
\usetheme{Boadilla}
\usecolortheme{seahorse}
\setbeamertemplate{navigation symbols}{}
\setbeamertemplate{itemize items}[circle]
\setbeamercolor{title}{bg=ucbblue, fg=white}
\setbeamercolor{frametitle}{bg=ucbblue, fg=white}

% --- Metadatos ---
\title[Estrategia Dual de Thévenin]{Estrategia Dual de Thévenin}
\subtitle{Del Sensor Pasivo a la Preamplificación Activa}
\author[M.Sc. Ing. B. Quiroga]{M.Sc. Ing. Bernardo Quiroga Turdera}
\institute[UCB Tarija]{IMT-121 Circuitos Electrónicos I | UCB Sede Tarija}
\date{Sesión 08}

\begin{document}

% ==========================================
% SLIDE 1: Objetivos y Mapa
% ==========================================
\begin{frame}
    \titlepage
\end{frame}

\begin{frame}{Resultados de Aprendizaje y Mapa de Navegación}
    \begin{block}{Objetivos de la Sesión}
        \begin{itemize}
            \item \textbf{Reducir} puentes de medición a modelos canónicos pasivos ($V_{th}, R_{th}$).
            \item \textbf{Demostrar cuantitativamente} la Máxima Transferencia de Potencia ($R_L = R_{th}$).
            \item \textbf{Aplicar} el método $V_{\text{test}} / I_{\text{test}}$ para caracterizar etapas activas lineales.
        \end{itemize}
    \end{block}
    
    \vspace{0.4cm}
    \begin{columns}
        \begin{column}{0.8\textwidth}
            \textbf{Estructura de la Clase:}
            \begin{enumerate}
                \item \textbf{Fase 1:} Dominio Pasivo (Sensor Desbalanceado).
                \item \textbf{Fase 2:} Extensión Activa (Preamplificación y Caja Negra).
            \end{enumerate}
        \end{column}
    \end{columns}
\end{frame}

% ==========================================
% SLIDE 2: El Puente de Medición
% ==========================================
\begin{frame}{El Puente de Medición (Wheatstone Pasivo)}
    El desbalance $\Delta R$ en un sensor produce una tensión diferencial $V_{oc} = V_B - V_A$.

    \begin{columns}
        \begin{column}{0.5\textwidth}
            \begin{center}
                \begin{circuitikz}[scale=0.75, transform shape]
                    \draw (0,2) node[left]{A} to[R, l=$R_3 {=} 350$] (2,0) node[below]{GND};
                    \draw (0,2) to[R, l=$R_1 {=} 350$] (2,4) node[above]{$V_s {=} \qty{10}{\volt}$};
                    \draw (2,4) to[R, l=$R_2 {=} 350$] (4,2) node[right]{B};
                    \draw (4,2) to[R, l=$R_4 {=} 360$] (2,0);
                    \draw (2,0) node[ground]{};
                \end{circuitikz}
            \end{center}
        \end{column}
        \begin{column}{0.5\textwidth}
            \textbf{Cálculo de $R_{th}$ (Apagando $V_s$ a GND):}
            \begin{center}
                \begin{circuitikz}[scale=0.6, transform shape]
                    \draw (0,0) node[left]{A} to[short, o-] (1,0) -- (1,1) to[R, l=$R_1$] (3,1) -- (3,0);
                    \draw (1,0) -- (1,-1) to[R, l=$R_3$] (3,-1) -- (3,0) node[ground]{};
                    \draw (3,0) -- (4,1) to[R, l=$R_2$] (6,1) -- (6,0);
                    \draw (3,0) -- (4,-1) to[R, l=$R_4$] (6,-1) -- (6,0);
                    \draw (6,0) to[short, -o] (7,0) node[right]{B};
                \end{circuitikz}
            \end{center}
            \begin{align*}
                R_A &= R_1 \parallel R_3 = 175\,\Omega \\
                R_B &= R_2 \parallel R_4 = 177.5\,\Omega \\
                R_{th} &= R_A + R_B = \mathbf{352.5\,\Omega}
            \end{align*}
        \end{column}
    \end{columns}
\end{frame}

% ==========================================
% SLIDE 3: Máxima Transferencia
% ==========================================
\begin{frame}{Máxima Transferencia de Potencia}
    La potencia transferida es máxima cuando la carga iguala la resistencia interna: $\mathbf{R_L = R_{th}}$.
    
    \vspace{0.4cm}
    \begin{columns}
        \begin{column}{0.4\textwidth}
            \begin{center}
                \begin{circuitikz}[scale=0.8, transform shape]
                    \draw (0,0) to[V, v=$V_{th}$] (0,2) 
                          to[R, l=$R_{th}$, -o] (3,2);
                    \draw (0,0) to[short, -o] (3,0);
                    \draw (3,2) to[vR, l=$R_L$, color=red, thick] (3,0);
                \end{circuitikz}
            \end{center}
        \end{column}
        \begin{column}{0.6\textwidth}
            \begin{alertblock}{Eficiencia Energética}
                A máxima transferencia, la eficiencia es de solo el \textbf{50\%} (la otra mitad se disipa térmicamente en $R_{th}$).
            \end{alertblock}
            \vspace{0.2cm}
            \textit{Justificación de Ingeniería:} En instrumentación biomédica y mecatrónica, priorizamos la extracción de la \textbf{señal y potencia útil} del sensor sobre la eficiencia energética global.
        \end{column}
    \end{columns}
\end{frame}

% ==========================================
% SLIDE 4: Colapso Activo
% ==========================================
\begin{frame}{El Colapso del Método Pasivo en Redes Activas}
    \begin{columns}
        \begin{column}{0.6\textwidth}
            \textbf{¿Por qué NO podemos apagar una fuente dependiente?}
            \begin{itemize}
                \item Una fuente dependiente \textit{no} es un componente autónomo; es una función matemática lineal de una variable física interna ($E_1 = k V_x$).
                \item Apagarla asumiendo $E_1 = 0$ viola las leyes constitutivas del circuito activo.
                \item Inyecta energía dinámica, modificando la impedancia aparente.
            \end{itemize}
        \end{column}
        \begin{column}{0.4\textwidth}
            \begin{center}
                \begin{circuitikz}[scale=0.8, transform shape]
                    \draw (0,2) to[cV, v=$E_1 {=} 50V_x$] (0,0);
                    \draw[red, thick] (-0.5,-0.5) -- (1,2.5);
                    \draw[red, thick] (-0.5,2.5) -- (1,-0.5);
                    \node[red, below] at (0,-0.5) {\textbf{¡Error Físico!}};
                \end{circuitikz}
            \end{center}
        \end{column}
    \end{columns}
\end{frame}

% ==========================================
% SLIDE 5: Método Vtest
% ==========================================
\begin{frame}{El Método Universal de la Fuente de Prueba ($V_{\text{test}} / I_{\text{test}}$)}
    \begin{enumerate}
        \item \textbf{Apagar} todas las fuentes \textit{independientes} ($V_s = 0 \to \text{Corto}$, $I_s = 0 \to \text{Abierto}$).
        \item \textbf{Dejar} las fuentes controladas/dependientes funcionando normalmente.
        \item \textbf{Conectar} una fuente $V_{\text{test}} = \qty{1.0}{\volt}$ en la salida y calcular la corriente inyectada $I_{\text{test}}$.
        \item $R_{\text{out}} = \frac{V_{\text{test}}}{I_{\text{test}}}$.
    \end{enumerate}

    \vspace{0.2cm}
    \begin{center}
        \begin{circuitikz}[scale=0.75, transform shape]
            \draw (0,0) node[ground]{} to[R, l=$R_{th} {=} 352.5$] (0,2) -- (1,2) node[above]{Nodo 1};
            \draw (1,2) to[R, l=$R_{in} {=} 1\text{k}$, v=$V_x$] (1,0) node[ground]{};
            \draw (1,2) to[R, l=$R_f {=} 100$] (5,2) node[above]{Nodo C};
            \draw (5,2) to[short, *-o] (6,2) node[right]{$V_{\text{test}} {=} \qty{1.0}{\volt}$};
            \draw (5,2) to[cV, v=$E_1 {=} 50V_x$] (5,0) node[ground]{};
            \draw[->, thick, red] (6,2.3) -- (5.2,2.3) node[midway, above]{$I_{\text{test}}$};
        \end{circuitikz}
    \end{center}
\end{frame}

% ==========================================
% SLIDE 6: Protocolo y Checklist
% ==========================================
\begin{frame}{Protocolo de Validación y Autoevaluación}
    \textbf{Protocolo de Validación en LTspice:}
    \begin{itemize}
        \item \textbf{Simulación Pasiva:} Directiva \texttt{.step param RL 10 1000 10} para graficar la campana de potencia en la carga.
        \item \textbf{Simulación Activa:} Directiva \texttt{.tf V(OUT) V\_sensor} para verificar la impedancia de salida en pequeña señal ($R_{\text{out}}$).
    \end{itemize}

    \vspace{0.2cm}
    \begin{block}{Checklist de Autoevaluación para Estudiantes}
        \begin{itemize}
            \item[$\square$] ¿Sé calcular $V_{th}$ en un puente desbalanceado usando divisores de tensión?
            \item[$\square$] ¿Puedo explicar por qué $R_{th}$ pasivo son dos paralelos en serie al apagar la fuente?
            \item[$\square$] ¿Entiendo por qué la eficiencia es del $50\%$ a máxima transferencia?
            \item[$\square$] ¿Sé formular ecuaciones nodales con una fuente de prueba $V_{\text{test}} = \qty{1}{\volt}$ y fuentes controladas activas?
        \end{itemize}
    \end{block}
\end{frame}

\end{document}
